\section{References}\label{sec:references}

Sources are grouped by the role they play. Only the editions this book quotes or relies on are
listed, and each entry says what was used. Scripture in English is the Douay-Rheims, cited by its
own chapters and verses under the usual names of the books; the readings of the Mass are cited in the numbering of the
Lectionary, which follows the New American Bible and numbers the psalms as in the Hebrew. A
Father's words quoted in the body with a short credit are found here under the Father's name, and
those the \textit{Catena Aurea} gives under Thomas Aquinas; where no translator is named for a
passage, the English is the book's own (Appendix~\ref{sec:scope}).

\subsection{Scripture}

\begin{itemize}[leftmargin=1.3em, itemsep=0.3em]
\item \textit{The Holy Bible}, Douay-Rheims version as revised by Richard Challoner, in the
electronic transcription of Project Gutenberg (eBook 1581), which does not name the printing it
follows. All Scripture quoted in English except the two canticles of Appendix~\ref{sec:prayers},
cited by the Douay's chapters and verses.
\item \textit{The Holy Bible translated from the Latin Vulgate}, Challoner's text (Baltimore: John
Murphy Company, 1914). The canticles of Zachary and of Mary, Luke 1:46--55 and 68--79, New
Testament pp.~66--67.
\item \latin{Biblia Sacra Vulgatae editionis} (the Clementine Vulgate), in the electronic
transcription of eBible.org (latVUC). Scripture quoted in Latin.
\item \textit{The Greek New Testament: SBL Edition}, ed. Michael W. Holmes, version 1.2, in its
electronic text. Luke 2:7 and 23:53, for the verbs of wrapping and laying.
\end{itemize}

\subsection{Liturgical books}

\begin{itemize}[leftmargin=1.3em, itemsep=0.3em]
\item \latin{Missale Romanum}, editio typica tertia (Vatican City: Typis Vaticanis, 2002), read
and not printed: the Collects of the Sundays of Advent and of 17--24 December; the Entrance
Antiphons of the Sundays, with the direction on the Gloria and the Creed; the two Prefaces of
Advent and the third Preface of the Nativity; the rubric that sets the Vigil Mass of the
Nativity on the evening of 24 December; the Communion rite of the Order of Mass,
nn.~124--146. In the same volume, the General Instruction (\latin{Institutio generalis Missalis
Romani}), nn.~30, 53, 54, 346 and 355; the Universal Norms on the Liturgical Year, n.~5; and
Paul VI's apostolic letter \latin{Mysterii Paschalis}.
\item The emended reprint of the third edition (2008), as announced in \latin{Notitiae} 44 (2008),
nos.~503--504, p.~367, with its additions to the General Roman Calendar, p.~369.
\item \textit{The Roman Missal}, English translation according to the third typical edition, for
use in the dioceses of the United States of America (New Jersey: Catholic Book Publishing, 2011).
Named and not quoted. Its General Instruction in the United States English (2011, as emended to
2021), nn.~48, 62 and 79, as published by the United States Conference of Catholic Bishops.
\item \textit{Lectionary for Mass for Use in the Dioceses of the United States}, second typical
edition (Confraternity of Christian Doctrine, 1970, 1986, 1997, 1998, 2001), never printed here. Its
numbers and citations are cited as the United States Conference of Catholic Bishops gives them
with the daily readings on its website: the Advent and Christmas Masses of the three cycles, and
the feast of St Andrew.
\item \latin{Ordo lectionum Missae}, editio typica altera (Vatican City: Libreria Editrice
Vaticana, 1981), read and not printed: the decrees of 1969 and 1981; the introduction, nn.~65, 66,
93 and 94; the readings of the weekdays of Advent, nn.~175--201, with the note before n.~193,
p.~110, and the Alleluia verses of 17--24 December.
\item \latin{Calendarium Romanum} (Vatican City: Typis Polyglottis Vaticanis, 1969): the decree,
p.~5; the Universal Norms on the Liturgical Year and the Calendar, nn.~5, 16\,b, 17, 39--42 and
45--47, pp.~12--17, and the Table of Liturgical Days, pp.~20--22; the commentary on the calendar,
pp.~60--61.
\item United States Conference of Catholic Bishops, \textit{Liturgical Calendar for the Dioceses
of the United States of America}, 2025, 2026 and 2027 (Washington, DC). The ranks of the December
saints and the transfer of the Immaculate Conception in 2024 (2025, pp.~7 and 11--12); Advent 2026
day by day (2026, pp.~45--47, among them Our Lady of Guadalupe, p.~46, and the Nativity, p.~47;
2027, pp.~7, 9 and 11--13); the Thirty-second Sunday in Ordinary Time, 8 November 2026 (2026,
p.~43).
\item The Liturgy of the Hours, no English edition of which is quoted. Its appointments for Advent
(the second readings of the Office of Readings, the \latin{Te Deum}, the antiphons at Evening
Prayer, and the Gospel canticles at Morning Prayer and Night Prayer, as on 1 December 2026) are
taken from the daily Office of the Association épiscopale liturgique pour les pays
francophones, on its website, for the Advents of 2021--2026, and from Felix Just, ``Office of
Readings: Non-Biblical Second Readings'', on catholic-resources.org, as updated 22 November 2026.
\item \latin{Institutio generalis de Liturgia Horarum}, in the English translation of 1974 for
Australia, England and Wales, and Ireland, \textit{General Introduction on the Liturgy of the
Hours}, as published by the Liturgy Office of the Bishops' Conference of England and Wales.
Nn.~37--38 and 50, on Lauds and Vespers and the Gospel canticles; nn.~55--59, on the Office of
Readings; nn.~84--91 and 138, on Compline; nn.~237--239, on memorials from 17 to 24 December; in
the book's own words.
\item \latin{Breviarium Romanum}, pars hiemalis, editio prima post typicam (Ratisbon, New York and
Cincinnati: Pustet, 1888). The Latin of the Office: the responsory \latin{Aspiciens a longe} of
the First Sunday of Advent, pp.~133--134; the \latin{Alma Redemptoris Mater}, p.~130; the great
O antiphons and their rubric, pp.~156--157; the antiphon \latin{O admirabile commercium},
p.~231; the general rubrics, title XXXI, on the \latin{Te Deum}, p.~xxxviii.
\item \textit{The Roman Breviary}, tr. John, Marquess of Bute, new edition, vol.~I, Winter
(Edinburgh and London: Blackwood, 1908). The English of the same: the first responsory of the
First Sunday of Advent, p.~214; the \latin{Alma Redemptoris Mater}, pp.~210--211; the O antiphons,
pp.~243--244; St Gregory the Great's eighth homily on the Gospels, the seventh lesson of
Christmas, p.~276; the first antiphons of Lauds of New Year's Day, p.~315; the note on Compline.
Bute prints the Office as it stood before 1960.
\item The same, vol.~II, Spring (1908). St Gregory the Great's twenty-ninth homily on the Gospels,
the eighth lesson of the Octave of the Ascension, pp.~488--489.
\item \latin{Missale Romanum}, editio typica (Vatican City: Typis Polyglottis Vaticanis, 1962). The
Introit of the Vigil of the Nativity, p.~15.
\item \latin{Missale Romanum}, editio quarta juxta typicam Vaticanam (Tours: Mame, 1922). The
Introits of the Third and Fourth Sundays of Advent, pp.~6 and 16.
\item \latin{Missale Romanum} (Ratisbon: Pustet, 1862). The prayer after Communion of the votive
Mass of our Lady in Advent, \latin{Gratiam tuam}.
\item F.~X. Lasance and F.~A. Walsh, \textit{The New Roman Missal in Latin and English}, revised
edition (New York: Benziger, 1945). The English of the Introit of the Fourth Sunday of Advent,
p.~130, in the printing of 1946; the Collect of St John of the Cross, 24 November, p.~1296.
\item \textit{The Roman Missal Translated into the English Language for the Use of the Laity}, first
revised edition (Philadelphia: Cummiskey, 1861). \latin{Domine, non sum dignus}, in Latin and
English; the translator is not named.
\item \textit{The Gelasian Sacramentary: Liber sacramentorum Romanae ecclesiae}, ed. H.~A. Wilson
(Oxford: Clarendon Press, 1894). The five Advent Masses, Book II, LXXX--LXXXIV, pp.~214--219,
with the prayers behind the Collects of the First Sunday of Advent and of 18, 21 and 24
December; \latin{Festinantes}, behind the Collect of the Second Sunday.
\item \textit{The Gregorian Sacramentary under Charles the Great}, ed. H.~A. Wilson, Henry Bradshaw
Society 49 (London, 1915). The Annunciation's prayer after Communion, \latin{Gratiam tuam}, p.~25;
the Advent prayers and stations of December, pp.~113--116.
\item \textit{Sacramentarium Leonianum}, ed. C.~L. Feltoe (Cambridge: University Press, 1896). An
Ascension preface, p.~21, and a December preface, pp.~162--163.
\item William Bright, \textit{Ancient Collects and Other Prayers}, 6th edition (Oxford and London:
Parker, 1887). The English of two Gelasian Advent prayers, pp.~16 and 18, and of a Veronese
Ascension collect, p.~60, whose clause on the second coming is the Ascension preface's.
\item \latin{Manuale Ambrosianum}, pars altera, ed. M. Magistretti (Milan: Hoepli, 1904). The six
Sundays of the Ambrosian Advent, pp.~1--41.
\item Anthony Ward, on the Rotulus of Ravenna, \latin{Notitiae} 25 (1989), no.~273, pp.~351--382.
The Rotulus prayers taken into the Advent Collects of the present Missal.
\item \latin{Martyrologium Romanum} (Rome: Propaganda Fide, 1902). St Cyril of Jerusalem,
18 March, p.~39; St Cyprian, 14 and 16 September, pp.~137--138; St John of the Cross,
24 November, p.~177, and 14 December, pp.~187--188; St Francis Xavier, 3 December, p.~182;
St Nicholas, 6 December, pp.~183--184; St Ambrose, 7 December, p.~184; St Damasus, 11 December,
p.~186; St Lucy, 13 December, p.~187.
\item \textit{The Little Office of the Blessed Virgin Mary}, Latin and English (New York: Benziger,
1915). Its rubrics, nos.~7, 8 and 14, pp.~12--13.
\end{itemize}

\subsection{Law and the Holy See}

\begin{itemize}[leftmargin=1.3em, itemsep=0.3em]
\item Council of Trent, Session XIV, chapters 3 and 4, on the parts of penance and on contrition:
\latin{Canones et decreta} (Leipzig: Tauchnitz, 1887), p.~73, and \textit{The Canons and Decrees
of the Sacred and Oecumenical Council of Trent}, tr. J. Waterworth (London: Burns and Oates,
undated reprint of the edition of 1848), p.~95.
\item Pius IX, \latin{Ineffabilis Deus} (1854), in \textit{Official Documents Connected with the
Definition of the Dogma of the Immaculate Conception}, Latin and English (Baltimore: Murphy;
London: Dolman, 1855). The Fathers on Genesis 3:15, pp.~77--79, and the definition, pp.~94--95.
The English translation is anonymous.
\item John Paul II, homily at the beatification of Pius IX, John XXIII and three others,
3 September 2000, English as published on the website of the Holy See. The heading and n.~1.
\item Second Vatican Council, constitution \latin{Sacrosanctum Concilium}, 4 December 1963,
\latin{Acta Apostolicae Sedis} 56 (1964). Nn.~51, 89, 100, 102 and 107, pp.~114--126, in the book's
own words.
\item Paul VI, apostolic constitution \latin{Missale Romanum}, 3 April 1969, \latin{Acta
Apostolicae Sedis} 61 (1969), pp.~217--222, in the book's own words.
\item Paul VI, apostolic constitution \latin{Laudis canticum}, 1 November 1970, \latin{Acta
Apostolicae Sedis} 63 (1971), pp.~527--535, with n.~2, pp.~529--530, in the book's own words.
\item Paul VI, apostolic exhortation \latin{Marialis cultus}, 2 February 1974, \latin{Acta
Apostolicae Sedis} 66 (1974). Nn.~3--4, pp.~117--118, and n.~41 with its note, p.~152, in the book's
own words.
\item \latin{Codex Iuris Canonici} (1983), Latin and English, as published on the website of the
Holy See. Canons 1, 1246--1247 and 1249--1253.
\item Congregation for Divine Worship and the Discipline of the Sacraments, \textit{Directory on
Popular Piety and the Liturgy} (2001), English as published on the website of the Holy See.
Nn.~98--105, in the book's own words.
\item Congregation for Divine Worship and the Discipline of the Sacraments, decree of 7 October
2019 on the memorial of Our Lady of Loreto, as published by the Holy See Press Office on
31 October 2019, in the book's own words.
\item Francis, apostolic letter \latin{Admirabile signum}, 1 December 2019, as published on the
website of the Holy See. Nn.~2, 3, 5 and 6 and the place and date of signing, in the book's own
words.
\item Dicastery for Divine Worship and the Discipline of the Sacraments, \textit{Note on the
Obligation to Attend Mass When a Holyday of Obligation Is Transferred}, 23 January 2025, official
English. Nn.~1--7.
\item National Conference of Catholic Bishops (United States), \textit{Pastoral Statement on
Penance and Abstinence}, 18 November 1966, as published by the United States Conference of
Catholic Bishops. Paragraphs 7, 17 and 22--23 and the note on its effective date.
\item United States Conference of Catholic Bishops, complementary norms on canon 1246 \S\,2 (the
decree on holy days of obligation, in force from 1 January 1993) and on canons 1252--1253, as
published by the Conference.
\item United States Conference of Catholic Bishops, Secretariat of Divine Worship,
\textit{Newsletter} 59 (November 2023), p.~44, and 61 (February 2025), pp.~5--6, on the
Immaculate Conception transferred in 2024 and on the Dicastery's note.
\item United States Conference of Catholic Bishops, news release of 20 August 2010 on the Roman
Missal, third edition, with the decree that set its use in the United States from the First
Sunday of Advent, 27 November 2011.
\end{itemize}

\subsection{The Fathers}

\begin{itemize}[leftmargin=1.3em, itemsep=0.3em]
\item Ambrose, \textit{Selected Works and Letters}, Nicene and Post-Nicene Fathers, 2nd series,
vol.~10, in the electronic transcription of the Christian Classics Ethereal Library. The editors'
note on his hymns.
\item Athanasius, \textit{On the Incarnation of the Word}, tr. A. Robertson, Nicene and
Post-Nicene Fathers, 2nd series, vol.~4 (New York, 1892). Chapter 54.2--3.
\item Augustine, \textit{Confessions}, tr. J.~G. Pilkington, Nicene and Post-Nicene Fathers, 1st
series, vol.~1 (Buffalo, 1886). VIII.12.29.
\item Augustine, \textit{The City of God}, Nicene and Post-Nicene Fathers, 1st series, vol.~2, in
the electronic transcription of the Christian Classics Ethereal Library. XVII.8, XVIII.30 and
XVIII.33.
\item Augustine, \textit{Sermons on Selected Lessons of the New Testament}, Nicene and Post-Nicene
Fathers, 1st series, vol.~6, in the same transcription. Sermons 12, 27 and 43 (62, 77 and 93 in
the Benedictine numbering).
\item Augustine, \textit{Expositions on the Book of Psalms}, Nicene and Post-Nicene Fathers, 1st
series, vol.~8 (New York, 1888), in the condensed English of A. Cleveland Coxe. On Psalm 37 (38),
13.
\item Augustine, Sermons 185 and 293, in \latin{Patrologia Latina}, vol.~38 (Paris: Migne, 1845),
cols.~997--998 and 1328--1329.
\item Cyprian, \textit{On the Advantage of Patience}, Ante-Nicene Fathers, vol.~5, in the
electronic transcription of the Christian Classics Ethereal Library. Sections 13 and 15.
\item Cyril of Jerusalem, \textit{Catechetical Lectures}, tr. E.~H. Gifford, Nicene and
Post-Nicene Fathers, 2nd series, vol.~7 (New York, 1894). Lecture XV, 1--3 and 26; the
Procatechesis, 9, in the electronic text of New Advent; and Gifford's introduction, chapter I, the
life of St Cyril, in the electronic transcription of the Christian Classics Ethereal Library.
\item \textit{The Epistle to Diognetus}, Ante-Nicene Fathers, vol.~1 (Buffalo, 1887). VIII.4 to
IX.6.
\item Gregory Nazianzen, \textit{Orations}, tr. C.~G. Browne and J.~E. Swallow, Nicene and
Post-Nicene Fathers, 2nd series, vol.~7 (New York, 1894). Oration 29, 19; Oration 38, 1 and 17;
Oration 45, 9.
\item Gregory the Great, \latin{Homiliae in Evangelia}, in \latin{Patrologia Latina}, vol.~76
(Paris: Garnier, 1878). Homilies 8.1, 12.6, 13.3 and 29.10, cols.~1103--1104, 1121--1122, 1124 and
1218--1219.
\item Hippolytus, \textit{Against the Heresy of One Noetus}, Ante-Nicene Fathers, vol.~5, in the
electronic transcription of the Christian Classics Ethereal Library. Sections 9 and 12.
\item Irenaeus, \textit{Against Heresies}, tr. A. Roberts and W.~H. Rambaut, Ante-Nicene Fathers,
vol.~1 (Buffalo, 1887). III.20.2--3, III.22.4, IV.20.4, the preface of Book V, and V.19.1.
\item Jerome, \latin{Commentarii in Isaiam}, Book IV, on Isaiah 11, in \latin{Patrologia Latina},
vol.~24 (Paris: Migne, 1865), cols.~147--148.
\item John Chrysostom, \textit{Homilies on the Gospel of Saint Matthew}, tr. G. Prevost, Nicene and
Post-Nicene Fathers, 1st series, vol.~10 (New York, 1888). Homilies 77.2--3, 78.1--2 and 79.1--2.
\item John Chrysostom, \textit{Homilies on the Gospel of St John}, tr. C. Marriott, Nicene and
Post-Nicene Fathers, 1st series, vol.~14 (New York, 1889). Homily 19.1, p.~67.
\item John of Damascus, \textit{An Exact Exposition of the Orthodox Faith}, tr. S.~D.~F. Salmond,
Nicene and Post-Nicene Fathers, 2nd series, vol.~9 (New York: Scribner, 1908). III.12, p.~56.
\item Leo the Great, \textit{Sermons} and \textit{Letters}, tr. C.~L. Feltoe, Nicene and
Post-Nicene Fathers, 2nd series, vol.~12 (New York, 1895). Sermons 19.2--3, 21.1 and 21.3, and
54.4; Letter 31.2--3.
\end{itemize}

\subsection{Saints and later writers}

\begin{itemize}[leftmargin=1.3em, itemsep=0.3em]
\item Alphonsus Liguori, \textit{The Incarnation, Birth and Infancy of Jesus Christ}, ed. E. Grimm,
Centenary Edition, vol.~IV, second edition (New York: Benziger, [1886]). The editor's notice,
p.~12, and the meditations for the novena of Christmas, 16--24 December, pp.~214--237.
\item Amalarius of Metz, \latin{Liber officialis} III.40, in \latin{Patrologia Latina}, vol.~105
(Paris: Migne), cols.~1158--1159, with the volume's table of contents, which places him in the
ninth century.
\item Anselm, \textit{Proslogium}, tr. S.~N. Deane (Chicago: Open Court, 1903). Chapter I, pp.~3
and 6.
\item Bernard of Clairvaux, \textit{St Bernard's Sermons for the Seasons and Principal Festivals
of the Year}, vol.~I, tr. a priest of Mount Melleray (Dublin: Browne and Nolan, [1921]). The fifth
sermon for Advent, 1 and 3, and the fourth sermon on the Virgin Mother (\latin{Super Missus est}
IV.8).
\item Bernard of Clairvaux, \textit{Sermons of St Bernard on Advent and Christmas}, tr. a member
of the community of St Mary's Convent, York (London: Washbourne, 1909). The first sermon on the
Vigil of the Nativity, 6.
\item Bernard of Clairvaux, \latin{Super Missus est} IV.8, in \latin{Patrologia Latina}, vol.~183
(Paris: Migne, 1854), col.~84, for the Latin quoted.
\item Bonaventure, \textit{The Life of Saint Francis}, tr. E. Gurney Salter (London: Dent, 1904).
X.7, pp.~110--112.
\item Francis Xavier, letter to his brethren at Rome from Cochin, 31 December 1543, in
H.~J. Coleridge, \textit{The Life and Letters of St Francis Xavier}, vol.~I (London: Burns and
Oates, 1872), pp.~155--156, in Coleridge's English; and Coleridge's own narrative,
pp.~287--288, on the voyages that led Xavier on to Japan.
\item Prosper Guéranger, \textit{The Liturgical Year: Advent}, tr. Laurence Shepherd, second
edition (Dublin: Duffy, 1870). On St Perpetuus, p.~24; on Durandus and the Advent fast in
France, p.~26; the \latin{Rorate caeli}, p.~155.
\item Honorius Augustodunensis, \latin{Gemma animae} III.1, in \latin{Patrologia Latina},
vol.~172 (Paris: Migne, 1854), cols.~641--643, and \latin{Sacramentarium} XCI, in the same volume
(Paris: Garnier, 1895), col.~797.
\item Ignatius of Loyola, \textit{The Spiritual Exercises}, tr. E. Mullan (New York: Kenedy, 1914).
The general examen, pp.~30--31.
\item Isaac of Stella, in \latin{Patrologia Latina}, vol.~194 (Paris: Migne). The head-title of
his works, naming him a Cistercian and abbot of Stella.
\item John Beleth, \latin{Summa de ecclesiasticis officiis} LXIV, in \latin{Patrologia Latina},
vol.~202 (Paris: Migne), col.~71.
\item John of the Cross, \textit{The Complete Works}, vol.~II, tr. David Lewis (London: Longman,
Green, 1864). Spiritual Maxims 70, p.~354; the romances \textit{The Desires of the Holy Fathers},
pp.~426--427, \textit{The Incarnation}, with the romance that follows it, p.~431, and
\textit{The Nativity}, p.~432.
\item Luis Laso de la Vega, \latin{Huei tlamahuiçoltica} (Mexico: Juan Ruiz, 1649). Our Lady's
words to Juan Diego, fol.~5r.
\item Thérèse of Lisieux, \textit{The Autobiography}, in \textit{Sœur Thérèse of Lisieux, the
Little Flower of Jesus}, ed. T.~N. Taylor (London: Burns and Oates, [1912]). Chapter V,
Christmas 1886, pp.~71--72. No translator is named.
\item Thomas Aquinas, \textit{Catena Aurea}, the Oxford translation, ed. J.~H. Newman: St Matthew,
tr. Mark Pattison (Oxford: Parker, 1842), on 8:8 and 25:9, and St Mark, tr. John Dalgairns
(Oxford: Parker, 1842), on 13:35 and 13:37, both in the electronic transcription of the Christian
Classics Ethereal Library; St Luke, tr. T.~D. Ryder (Oxford: Parker, 1843), on 1:39, 1:46--47
and 2:6--7, pp.~67--68. The Fathers as the Catena gives them: Ambrose, Augustine (Letter 199),
Bede, Chrysostom, Cyril, Gregory, Hilary, an unknown author under Origen's name, and Theophylact.
\item William Durandus, \latin{Rationale divinorum officiorum}, tomus secundus (Lyon, 1612).
VI.2, on the Gloria in Advent. And, from the edition of [1568] that binds it with John Beleth's
\latin{Rationale}, the head-title naming him bishop of Mende.
\end{itemize}

\subsection{History and councils}

\begin{itemize}[leftmargin=1.3em, itemsep=0.3em]
\item The calendar of 354: \latin{Feriale Ecclesiae Romanae}, in \latin{Chronographus anni
CCCLIIII}, ed. Th. Mommsen, Monumenta Germaniae Historica, Auctores antiquissimi 9 (Berlin:
Weidmann, 1892), p.~71. The entry for 25 December.
\item Karl Joseph von Hefele, \textit{A History of the Councils of the Church}, vol.~II, tr.
H.~N. Oxenham (Edinburgh: T. and T. Clark, 1876), pp.~292--293, the council of Saragossa, canon 4;
vol.~IV, tr. W.~R. Clark (Edinburgh: T. and T. Clark, 1895), pp.~403--404, the council of Mâcon,
canon 9.
\item M.~B. Riddle, introductory essay to Augustine's \textit{Harmony of the Gospels}, Nicene
and Post-Nicene Fathers, 1st series, vol.~6, in the electronic transcription of the Christian
Classics Ethereal Library. Its account of Tatian's \textit{Diatessaron}, the earliest harmony of
the four Gospels.
\item \latin{Concilia aevi Merovingici}, ed. F. Maassen, Monumenta Germaniae Historica, Concilia I
(Hanover: Hahn, 1893). The council of Mâcon, canon 9, for the Latin quoted.
\item Gregory of Tours, \latin{Historia Francorum} X.31, ed. W. Arndt, Monumenta Germaniae
Historica, Scriptores rerum Merovingicarum I (Hanover: Hahn, 1884), and \textit{The History of the
Franks}, tr. O.~M. Dalton, vol.~II (Oxford: Clarendon Press, 1927), pp.~471--472. The fasts of
St Perpetuus.
\item \latin{Concilia provincialia Baltimori habita ab anno 1829 usque ad annum 1849}, second
edition (Baltimore: Murphy, 1851). The Sixth Provincial Council of Baltimore, May 1846: its
election of the Blessed Virgin Mary conceived without stain as patroness of the United States.
\end{itemize}

\subsection{Prayer books and hymns}

\begin{itemize}[leftmargin=1.3em, itemsep=0.3em]
\item \textit{A Manual of Prayers for the Use of the Catholic Laity} (New York: Catholic Publication
Society, 1889). The act of contrition, p.~54.
\item \textit{The Raccolta}, tr. Ambrose St John (London: Burns and Oates, 1910). The Angelus,
no.~182, in Latin and English, pp.~151--152.
\item A.~A. Lings, \textit{Our Favorite Novenas} (New York: Benziger, 1897). The prayer said from
St Andrew's day to Christmas, p.~29.
\item \textit{A Manual of Gregorian Chant} (Rome and Tournai: Desclée, 1903). The \latin{Rorate
caeli} with its verses, pp.~313--315.
\item J.~H. Newman, \latin{Hymni ecclesiae} (Oxford: Parker, 1838). \latin{Veni, Redemptor
gentium}, pp.~89--90.
\item \textit{Hymnal Noted}, parts I and II (London: Novello, Ewer; Masters, 1851). ``Come, thou
Redeemer of the earth'', hymn 12, stanzas 1, 4, 5 and 7, pp.~35--36. The translator is not named
there; the Nicene and Post-Nicene Fathers call the version Dr Neale's.
\end{itemize}
